\documentclass[a4paper]{article}

\usepackage{graphicx}
\usepackage{amsmath,amssymb}
\usepackage{minted}
\usepackage{multirow}
\usepackage{float}
\usepackage{algorithm}
\usepackage{algpseudocode}
\usepackage{todonotes}
\usepackage{forest}
\usepackage{xstring}
\usepackage{xcolor}
\usepackage[most]{tcolorbox}
\usepackage{xparse}
\usepackage{natbib}
\usepackage{adjustbox}

\usetikzlibrary{graphs,graphdrawing}
\usegdlibrary{layered}

% for external graphviz
\input{/home/donkarlo/Dropbox/repo/nd_language_ai_project/src/nd_language_ai/formal/computer/programming/tex/latex/code_clips/external_graphviz.tex}

% for tags
\input{/home/donkarlo/Dropbox/repo/nd_language_ai_project/src/nd_language_ai/formal/computer/programming/tex/latex/part/control_sequence/kind/semtag/semtag.tex}

% for links
\input{/home/donkarlo/Dropbox/repo/nd_language_ai_project/src/nd_language_ai/formal/computer/programming/tex/latex/code_clips/seehere.tex}

% for nested lists and sections
\input{/home/donkarlo/Dropbox/repo/nd_language_ai_project/src/nd_language_ai/formal/computer/programming/tex/latex/code_clips/deep_structure.tex}

\title{Prediction Residuals for Outlier Detection}
\date{}
\author{Mohammad Rahmani}

\begin{document}
   \maketitle

\section{Scoring rule}
Use a model to estimate the expected value from context. For a time series, form a one-step-ahead prediction \(\hat y_{t\mid t-1}\) using data available before time \(t\). The residual is
\[
e_t=y_t-\hat y_{t\mid t-1}.
\]
If predictive uncertainty varies, score \(|e_t|/\hat\sigma_t\) with \(\hat\sigma_t>0\), or flag observations outside a calibrated prediction interval.

\section{Example}
A service normally expects \(100\) requests per minute at a particular hour, with a predictive standard deviation of \(10\). A count of \(140\) gives a standardized residual of \(4\). The same count could be ordinary during a busier hour.

\section{Practical considerations}
Use chronological validation to prevent leakage and account for seasonality and changing variance. Persistent large residuals can indicate a regime change or a poor model. In a Kalman filter, the measurement minus its prior prediction is the innovation and can be scaled by its predicted covariance.



   \bibliographystyle{plainnat}
   \bibliography{/home/donkarlo/Dropbox/repo/nd_shared_research/bibliography/bibliography.bib}
\end{document}
